

\medskip
\noindent\SGARefTarget{sga3:iv:lemma:4:4:7}\textbf{Lemma 4.4.7.}
\textit{Every family of morphisms of sheaves \(\{F_i\to F\}\) factors as a
covering family \(\{F_i\to G\}\) followed by a monomorphism \(G\to F\).}

Indeed, it is enough to take \(G\) to be the sheaf image of the given family.

\medskip
\noindent\textit{Proof of Proposition \SGARefLink{sga3:iv:proposition:4:4:3}{4.4.3}.}
We saw in \SGARefLink{sga3:iv:lemma:4:4:6}{4.4.6} that \SGARefLink{sga3:iv:proposition:4:4:3:item:iii:01}{(iii)} implies \SGARefLink{sga3:iv:proposition:4:4:3:item:ii:01}{(ii)}, while \SGARefLink{sga3:iv:proposition:4:4:3:item:ii:01}{(ii)} evidently implies \SGARefLink{sga3:iv:proposition:4:4:3:item:i:01}{(i)}.
Finally, let \(\{F_i\to F\}\) be an epimorphic family. By Lemma~\SGARefLink{sga3:iv:lemma:4:4:7}{4.4.7} it
factors as a covering family followed by a monomorphism. The latter
monomorphism, being dominated
by an epimorphic family, is itself an epimorphism; hence it is an isomorphism
by \SGARefLink{sga3:iv:lemma:4:4:4}{4.4.4}.

\medskip
\noindent\SGARefTarget{sga3:iv:remark:4:4:8}\textbf{Remark 4.4.8.}
Because the presheaf image of the family
\(\mathscr F=\{F_i\to F\}\) is separated, the construction of the associated
sheaf shows that the conditions of Proposition~\SGARefLink{sga3:iv:proposition:4:4:3}{4.4.3} are also equivalent to
the following:
\begin{enumerate}
  \item[(v)]\SGARefTarget{sga3:iv:remark:4:4:8:item:v:01} for every \(S\in\operatorname{Ob}\mathscr C\), each
    \(f\in F(S)\) lies locally in the image of \(\mathscr F\); that is,

  \item[(vi)]\SGARefTarget{sga3:iv:remark:4:4:8:item:vi:01} for every \(S\in\operatorname{Ob}\mathscr C\) and every
    \(f\in F(S)\), the set of morphisms \(S'\to S\) such that the image of
    \(f\) in \(F(S')\) belongs to the image of one of the \(F_i(S')\) is a
    refinement of \(S\);

  \item[(vii)]\SGARefTarget{sga3:iv:remark:4:4:8:item:vii:01} if the topology is defined by a pretopology
    \(S\mapsto\mathscr R(S)\), then, for every
    \(S\in\operatorname{Ob}\mathscr C\) and every \(f\in F(S)\), there is a
    family \(\{S_j\to S\}\in\mathscr R(S)\) such that, for each \(j\), the
    image \(f_j\) of \(f\) in \(F(S_j)\) belongs to the image of one of the
    \(F_i(S_j)\).
\end{enumerate}

\medskip
\noindent\SGARefTarget{sga3:iv:remark:4:4:8:bis}\textbf{Remark 4.4.8 bis.}
If the sheaf \(F\) is representable, the preceding conditions are also
equivalent to:
\begin{enumerate}
  \item[(viii)]\SGARefTarget{sga3:iv:remark:4:4:8:bis:item:viii:01} the set of morphisms \(T\to F\), with
    \(T\in\operatorname{Ob}\mathscr C\), for which there exist an \(i\) and
    a commutative diagram
    \[
    \begin{array}{ccc}
       &F_i&\\[-0.4ex]
       &\big\downarrow&\\[-0.4ex]
      T&\longrightarrow&F
    \end{array}
    \]
    is a refinement of \(F\).
\end{enumerate}

Indeed, if \SGARefLink{sga3:iv:remark:4:4:8:bis:item:viii:01}{(viii)} holds, the presheaf image of the \(F_i\) is dominated
by
a refinement of \(F\), which implies that the family is covering.
Conversely, apply \SGARefLink{sga3:iv:remark:4:4:8:item:vi:01}{(vi)} to \(\operatorname{id}_F\in F(F)\).

\SGARefTarget{sga3:iv:remark:4:4:8:bis:item:b:01}In pictorial language, \SGARefLink{sga3:iv:remark:4:4:8:bis:item:viii:01}{this condition} says: \textit{locally on \(F\), there
is an \(i\) for which \(F_i\to F\) has a section}. In particular, a morphism
\[
  G\longrightarrow F,
\]
where \(G\) is a sheaf and \(F\) is a representable sheaf, is covering if
and only if it has a section locally on \(F\).
